\scanpage{226}
% archive sheet number: 225

\mspage{- 6 -} \hfill n\textsuperscript{o}370

2. \uline{Le théorème d'existence et d'unicité}.

\figure{\begin{tikzpicture}[x=1cm,y=1cm,gray!55,line width=.27pt]\draw (0,.12)--(.11,-.03) (.07,.14)--(.17,-.04) (.14,.15)--(.24,0) (.13,.05)--(.25,.18) (.21,.12)--(.31,.03);\end{tikzpicture}}

DÉFINITION 2.- \uline{Soit} $G$\handwritten{\unsure{$\checkmark$}} \uline{un groupe localement compact}. \uline{On appelle mesure de Haar à gauche} (resp. \uline{à droite}) \uline{sur} $G$ \uline{une mesure positive non nulle sur} $G$, \uline{invariante à gauche} (resp. \uline{à droite}).

THÉORÈME 1.- \uline{Sur tout groupe localement compact}, \uline{il existe une mesure de Haar à gauche} (resp. \uline{à droite}), \uline{et}, \uline{à un facteur constant près}, \uline{il n'en existe qu'une}.

A) \uline{Existence} - Posons $\mathcal{K}(G)=\mathcal{K}$, $\mathcal{K}_+(G)=\mathcal{K}_+$, $\mathcal{K}_+^*=\mathcal{K}_+-\{0\}$. Si $C$ est une partie compacte de $G$, on notera $\mathcal{K}_+^*(C)$ l'ensemble des $f\in\mathcal{K}_+^*$ à support dans $C$. Pour $f\in\mathcal{K}$ et $g\in\mathcal{K}_+^*$, il existe des nombres $c_1,\ldots,c_n\geq0$ et des éléments $s_1,\ldots,s_n$ de $G$ tels que $f\leq\sum\limits_{i=1}^n c_i\,\gamma(s_i)g$. (En effet, il existe une partie ouverte non vide $U$ de $G$ telle que $\inf_{s\in U}g(s)>0$, et le support de $f$ peut être recouvert par un nombre fini de translatés à gauche de $U$). Soit alors $(f:g)$ la borne inférieure des nombres $\sum\limits_{i=1}^n c_i$ pour tous les systèmes $(c_1,\ldots,c_n,s_1,\ldots,s_n)$ de nombres $\geq0$ et d'éléments de $G$ tels que $f\leq\sum\limits_{i=1}^n c_i\,\gamma(s_i)g$. On a :

\[
(\gamma(s)f:g)=(f:g)\qquad\text{pour }f\in\mathcal{K},\;g\in\mathcal{K}_+^*,\;s\in G\ ; \tag*{(i)}
\]

\[
(\lambda f:g)=\lambda(f:g)\qquad\text{pour }f\in\mathcal{K},\;g\in\mathcal{K}_+^*,\;\lambda\geq0\ ; \tag*{(ii)}
\]

\[
(f+f':g)\leq(f:g)+(f':g)\qquad\text{pour }f,f'\in\mathcal{K},\;g\in\mathcal{K}_+^*\ ; \tag*{(iii)}
\]

\[
(f:g)\geq(\sup f)/(\sup g)\qquad\text{pour }f\in\mathcal{K},\;g\in\mathcal{K}_+^*\ ; \tag*{(iv)}
\]

\[
(f:h)\leq(f:g)(g:h)\qquad\text{pour }f\in\mathcal{K},\;g,h\in\mathcal{K}_+^*\ ; \tag*{(v)}
\]

\[
0<\frac{1}{(f_{\unsure{o|0}}:f)}\leq\frac{(f:g)}{(f_o:g)}\leq(f:f_o)\qquad\text{pour }f,f_o,g\in\mathcal{K}_+^*\ ; \tag*{(vi)}
\]

\[
\begin{aligned}
&\text{soient } f,f',h\in\mathcal{K}_+\text{ avec }h(s)\geq1\text{ sur le support de }f+f',\text{ et soit}\\
&\varepsilon>0\ ;\ \text{il existe un voisinage compact }V\text{ de }e\text{ tel que, pour toute}\\
&g\in\mathcal{K}_+^*(V),\ \text{on ait }(f:g)+(f':g)\leq(f+f':g)+\varepsilon(h:g).
\end{aligned} \tag*{(vii)}
\]

Les propriétés (i), (ii), (iii) sont évidentes. Soient $f\in\mathcal{K}$, $g\in\mathcal{K}_+^*$ ;
